%%%%%%%%%%%%%%%%%%%%%%%%%%%%%%%%%%%%%%%%%%%%%%%%%%%%%%%%%%%%%
% ==================== 第二章：问题分析 ====================
% BZD 数模社制作——国赛 B 题四问通用模板
% 使用时请根据题目和附件逐一替换【】中的内容，并删除不适用的表述。
%
% ✓ 自查清单：
%   □ 是否说明了每个问题的任务、研究意义和核心困难？
%   □ 是否判断了各问题的数学类型，并说明方法选择依据？
%   □ 是否分析了附件数据的内容、结构、质量和适用范围？
%   □ 是否说明了数据分析、清洗和预处理方法及其理由？
%   □ 是否明确了题目要求的输出、约束和评价标准？
%   □ 是否仅在确有必要时设置主模型与对比模型？
%   □ 是否说明了四个问题之间的数据或逻辑联系？
%   □ 是否规划了误差分析、灵敏度分析或稳健性检验？
%%%%%%%%%%%%%%%%%%%%%%%%%%%%%%%%%%%%%%%%%%%%%%%%%%%%%%%%%%%%%

\section{问题分析}

\subsection{问题一的分析}

问题一要求在【研究背景或已知条件】下完成【问题一的具体任务】，其研究意义在于【说明该结果对认识研究对象、揭示基本规律或解决实际问题的作用】。本问需要重点回答【核心问题】，并为后续问题提供【参数、特征、状态量、中间结果或基础结论】。

从数学建模角度看，问题一属于【机理分析、统计分析、预测、分类、优化或其他数学问题类型】，核心在于【需要刻画的关系、规律或决策目标】。解决此类问题通常可以采用【候选方法类别】，但不同方法对【样本规模、变量关系、分布特征、约束条件或先验假设】的要求不同，因此需要结合题目条件和附件数据进行选择，而不能仅依据模型复杂程度确定方案。

附件【附件名称或编号】主要包含【数据对象、变量和时间或空间范围】。应首先通过【描述性统计、分布分析、相关性分析、时序或空间特征分析】认识数据，并检查【缺失值、异常值、重复记录、量纲差异、统计口径或变量相关性】。针对发现的问题，采用【缺失值处理、异常值识别、数据变换、标准化或特征构造方法】，其理由是【说明该处理对保持数据真实性、满足模型条件或提高结果可靠性的作用】。

题目要求输出【问题一的结果形式】，并满足【精度、范围、物理意义或其他约束】，可利用【误差指标、拟合优度、约束回代或其他评价标准】检验结果。基于上述分析，首先建立【主模型名称或类型】刻画【主要关系】，必要时再建立【对比或备选模型】进行验证，并依据【比较指标】选择更合适的结果。问题一得到的【关键输出】将作为问题二中【对应变量或条件】的输入。

\subsection{问题二的分析}

问题二在问题一的【关键结果或基础模型】上进一步加入【新增条件、影响因素或研究范围】，要求完成【问题二的具体任务】。研究该问题有助于【说明其对模型深化、实际应用或决策改进的意义】，其核心是分析【新增因素】如何影响【研究对象或目标结果】。

从数学建模角度看，本问属于【参数估计、关系分析、预测、优化、评价或其他数学问题类型】，与问题一相比，新增难点在于【非线性、动态性、耦合关系、多重约束或不确定性】。因此，需要在问题一方法的基础上判断【原模型是否仍然适用】，并根据新增条件选择【模型扩展、参数修正、分层建模、组合建模或其他处理思路】。

本问使用的数据由【问题一的中间结果】与附件【附件名称或编号】中的【新增数据】共同构成。建模前需要检查二者在【时间范围、空间尺度、量纲、统计口径和对象标识】上的一致性，并通过【数据匹配、口径统一、相关性或因果性分析、特征筛选】识别真正影响本问结果的因素。若数据存在【样本不平衡、信息缺失、共线性或噪声干扰】，则采用【相应处理方法】，以避免【说明可能产生的偏差或失真】。

题目要求给出【问题二的结果形式】，并重点考察【准确性、稳定性、可解释性、计算效率或可实施性】。据此，可在问题一模型基础上构建【改进后的主模型】，并在有比较价值时选用【备选模型或基准方法】，利用【评价指标、交叉验证、对比实验或约束检验】比较结果。最终得到的【参数、预测结果、判定结果或方案】将为问题三中的【综合分析或最终决策任务】提供依据。

\subsection{问题三的分析}

问题三综合问题一和问题二得到的【基础结果与改进结果】，在【新的情景、约束或应用要求】下完成【问题三的具体任务】。本问是前述研究的进一步拓展，其意义在于【说明该结果对刻画复杂规律、完善模型或支持后续综合决策的价值】。

从数学建模角度看，本问属于【综合决策、多目标优化、情景分析、机理推演或其他数学问题类型】，核心难点是同时处理【多个目标、约束耦合、不确定因素或结果冲突】。因此，需要明确【决策变量或待求量】、【目标或评价准则】及【主要约束】，并分析问题一、问题二的结果如何转化为本问的【模型参数、边界条件或评价依据】。

本问需要整合附件【附件名称或编号】及前两问生成的【中间数据或模型输出】。应检查数据传递过程中的【误差累积、单位一致性、时间或空间对齐以及适用范围】，并通过【数据融合、情景划分、参数标定或不确定性描述方法】形成可用于求解的统一数据基础。对于附件无法直接给出的【参数或条件】，只能依据【明确的数据来源、题设条件或可验证的计算关系】确定，不得为了得到理想结果而人为设定。

题目要求输出【问题三的预测值、判定结果、优化方案或中间结论】，并满足【可行性、稳定性、精度或实际约束】。基于上述分析，可建立【综合主模型】完成求解；若确有必要，再以【基准模型、替代算法或不同情景】进行对比，并通过【误差分析、约束回代或对比实验】检验结果。问题三得到的【关键参数、情景结果或候选方案】将作为问题四中【综合决策或推广任务】的基础。

\subsection{问题四的分析}

问题四在前三问的【模型、参数、规律或候选方案】基础上，进一步考虑【更复杂的实际情景、新增限制或推广要求】，要求完成【问题四的具体任务】。本问承担全文的综合应用任务，其研究意义在于【说明最终结果对实际决策、系统改进、方案推广或管理建议的价值】。

从数学建模角度看，本问属于【综合决策、多目标优化、系统评价、情景推演或其他数学问题类型】，核心难点是协调【多个研究对象、多个目标、耦合约束、动态变化或不确定因素】。因此，需要明确【最终决策变量或待求量】、【总体目标或评价准则】和【必须满足的约束】，并说明前三问的结果如何转化为本问的【输入参数、边界条件、候选集合或评价依据】。

本问需要整合全部附件数据以及前三问生成的【中间数据、模型参数和求解结果】。应重点检查不同来源数据之间的【对象对应关系、时间或空间尺度、单位与统计口径】，并评估前序结果传递产生的【误差累积和不确定性】。针对【不同情景、参数波动或数据扰动】，可采用【情景划分、数据融合、不确定性描述或参数扰动方法】构造可比较的输入条件，从而保证最终结论具有明确的数据依据。

题目最终要求输出【问题四的综合结果、最优方案、实施策略或建议】，并满足【可行性、稳定性、经济性、时效性或其他实际约束】。基于上述分析，可建立【最终主模型】统一求解；若存在多种可行方法，则选用【基准模型、替代算法或不同情景】进行比较，并依据【综合评价指标】确定最终结果。最后通过【灵敏度分析、稳健性检验、误差分析、约束回代或对比实验】说明结论在【关键参数变化、数据扰动或应用场景改变】下是否仍然成立，并据此给出【与题目要求一致的最终结论或建议】。
